\section*{Chapter \ref{ch:g10:cellular-respiration} --- Cellular Respiration}

\begin{solution}{exo:g10:cellular-respiration:1}
The balanced equation is
$\mathrm{C_6H_{12}O_6}+6\,\mathrm{O_2}\to 6\,\mathrm{CO_2}+6\,\mathrm{H_2O}$
plus energy. Part of that energy is captured as \omterm{def:g10:cellular-respiration:ATP}{ATP}; the rest is released as
heat. The equation summarises aerobic oxidation of glucose, not the many
\omterm{def:g10:membranes-enzymes:enzyme}{enzyme} steps inside the \omterm{def:g10:the-cell:cell}{cell}.
\end{solution}

\begin{solution}{exo:g10:cellular-respiration:2}
\omterm{def:g10:cellular-respiration:glycolysis}{Glycolysis} occurs in the cytosol. The link reaction and \omterm{def:g10:cellular-respiration:krebs}{Krebs cycle} occur in
the mitochondrial \omterm{def:g10:cellular-respiration:mito}{matrix}. The \omterm{def:g10:cellular-respiration:ETC}{electron transport chain} and \omterm{def:g10:cellular-respiration:ETC}{oxidative
phosphorylation} are located on the inner mitochondrial membrane (\omterm{def:g10:cellular-respiration:mito}{cristae}).
Together these stages oxidise glucose fully when oxygen is available.
\end{solution}

\begin{solution}{exo:g10:cellular-respiration:3}
Oxygen is the final electron acceptor of the \omterm{def:g10:cellular-respiration:ETC}{electron transport chain}. It
combines with electrons and protons to form water. Without oxygen the chain
backs up, NADH cannot be reoxidised by the ETC, and \omterm{def:g10:cellular-respiration:ETC}{oxidative phosphorylation}
--- the main \omterm{def:g10:cellular-respiration:ATP}{ATP} source --- stops.
\end{solution}

\begin{solution}{exo:g10:cellular-respiration:4}
\omterm{def:g10:cellular-respiration:ATP}{ATP} is made and spent continuously for immediate cellular work (pumps,
synthesis, movement). It is not stockpiled in large amounts. Long-term energy
storage uses \omterm{def:g10:molecules-of-life:lipids}{fats}, starch or glycogen, which are mobilised and respired when
the \omterm{def:g10:the-cell:cell}{cell} needs to regenerate \omterm{def:g10:cellular-respiration:ATP}{ATP}.
\end{solution}

\begin{solution}{exo:g10:cellular-respiration:5}
\omterm{def:g10:cellular-respiration:def}{Aerobic respiration} requires oxygen, yields far more \omterm{def:g10:cellular-respiration:ATP}{ATP} per glucose, and
produces $\mathrm{CO_2}$ and $\mathrm{H_2O}$. \omterm{def:g10:cellular-respiration:ferment}{Fermentation} needs no oxygen,
yields only glycolytic \omterm{def:g10:cellular-respiration:ATP}{ATP}, and produces \omterm{def:g10:cellular-respiration:lactic}{lactate} (lactic \omterm{def:g10:cellular-respiration:ferment}{fermentation}) or
\omterm{def:g10:cellular-respiration:alcoholic}{ethanol} plus $\mathrm{CO_2}$ (\omterm{def:g10:cellular-respiration:alcoholic}{alcoholic fermentation}), depending on the
pathway.
\end{solution}

\begin{solution}{exo:g10:cellular-respiration:6}
Blocking the ETC stops \omterm{def:g10:cellular-respiration:ETC}{oxidative phosphorylation}, which supplies most \omterm{def:g10:cellular-respiration:ATP}{ATP} in
aerobic \omterm{def:g10:the-cell:cell}{cells}. Glucose may still be present, but it cannot be fully oxidised
for energy. Membrane pumps and other ATP-dependent processes fail, so tissues
--- especially brain and heart --- are damaged quickly.
\end{solution}

\begin{solution}{exo:g10:cellular-respiration:7}
\omterm{def:g10:cellular-respiration:glycolysis}{Glycolysis} requires NAD$^+$ to accept electrons when glucose is oxidised to
pyruvate. If the ETC is too slow to reoxidise NADH, \omterm{def:g10:cellular-respiration:ferment}{fermentation} transfers
electrons from NADH onto pyruvate (forming \omterm{def:g10:cellular-respiration:lactic}{lactate}), regenerating NAD$^+$ so
\omterm{def:g10:cellular-respiration:glycolysis}{glycolysis} --- and its small \omterm{def:g10:cellular-respiration:ATP}{ATP} yield --- can continue.
\end{solution}

\begin{solution}{exo:g10:cellular-respiration:8}
\omterm{def:g10:cellular-respiration:mito}{Cristae} fold the inner mitochondrial membrane, increasing \omterm{def:g10:the-cell:sa-v}{surface area} for
ETC complexes and \omterm{def:g10:cellular-respiration:ATP}{ATP} synthase without a huge increase in \omterm{def:g10:the-cell:prok-euk}{organelle} \omterm{def:g10:the-cell:sa-v}{volume}.
That is the same surface-to-volume logic as microvilli: more membrane work per
unit \omterm{def:g10:the-cell:sa-v}{volume}.
\end{solution}

\begin{solution}{exo:g10:cellular-respiration:9}
Stored food reserves are oxidised by \omterm{def:g10:cellular-respiration:def}{cellular respiration} to supply \omterm{def:g10:cellular-respiration:ATP}{ATP} for
germination. Carbon leaves the seed as $\mathrm{CO_2}$, so dry mass falls. In
darkness there is no \omterm{def:g10:photosynthesis:def}{photosynthesis} to replace organic mass, so the seed
continues to lose dry weight until light-driven carbon fixation begins.
\end{solution}

\begin{solution}{exo:g10:cellular-respiration:10}
\omterm{def:g10:scientific-biology:variables}{Independent variable}: temperature. \omterm{def:g10:scientific-biology:variables}{Dependent variable}: e.g.\ rate of manometer
fluid movement or change in a $\mathrm{CO_2}$ indicator. \omterm{def:g10:scientific-biology:variables}{Controlled variables}:
seed mass and type, apparatus \omterm{def:g10:the-cell:sa-v}{volume}, time, and KOH if measuring $\mathrm{O_2}$
uptake. Control: dead or boiled seeds, or glass beads, to show that abiotic
\omterm{def:g10:the-cell:sa-v}{volume} or indicator change is negligible.
\end{solution}

\begin{solution}{exo:g10:cellular-respiration:11}
$\mathrm{CO_2}$ is fixed in the \omterm{def:g10:photosynthesis:calvin}{Calvin cycle} into sugars and then stored as
starch in the leaf. An animal digests starch to glucose, absorbs it, and
respires it in mitochondria, exhaling $\mathrm{CO_2}$. Key processes:
\omterm{def:g10:photosynthesis:def}{photosynthesis}, digestion and absorption, \omterm{def:g10:cellular-respiration:def}{cellular respiration}, and
ventilation.
\end{solution}

\begin{solution}{exo:g10:cellular-respiration:12}
Lactic \omterm{def:g10:cellular-respiration:ferment}{fermentation} (muscle, some bacteria) converts pyruvate to \omterm{def:g10:cellular-respiration:lactic}{lactate}.
\omterm{def:g10:cellular-respiration:alcoholic}{Alcoholic fermentation} (yeast) converts pyruvate to \omterm{def:g10:cellular-respiration:alcoholic}{ethanol} and $\mathrm{CO_2}$.
Both regenerate NAD$^+$ so \omterm{def:g10:cellular-respiration:glycolysis}{glycolysis} can continue without oxygen. \omterm{def:g10:cellular-respiration:ATP}{ATP} yield
is low because only \omterm{def:g10:cellular-respiration:glycolysis}{glycolysis} produces \omterm{def:g10:cellular-respiration:ATP}{ATP}; the organic end products still
contain most of the energy that \omterm{def:g10:cellular-respiration:def}{aerobic respiration} would extract via Krebs
and the ETC.
\end{solution}

\begin{solution}{exo:g10:cellular-respiration:13}
\omterm{def:g10:cellular-respiration:lactic}{Lactate} is not made ``using oxygen''; it is made when oxygen delivery is
insufficient for the ETC. \omterm{def:g10:cellular-respiration:ferment}{Fermentation} regenerates NAD$^+$ so \omterm{def:g10:cellular-respiration:glycolysis}{glycolysis} can
keep making a little \omterm{def:g10:cellular-respiration:ATP}{ATP}. After the sprint, elevated breathing helps repay an
oxygen debt: restoring myoglobin/oxygen stores, supporting aerobic \omterm{def:g10:cellular-respiration:ATP}{ATP}
production, and allowing \omterm{def:g10:cellular-respiration:lactic}{lactate} clearance and metabolic recovery --- not
literally reversing \omterm{def:g10:cellular-respiration:lactic}{lactate} formation molecule-for-molecule with each extra
breath.
\end{solution}

\begin{solution}{pb:g10:cellular-respiration:1}
\textbf{Part I.} \omterm{def:g10:cellular-respiration:glycolysis}{Glycolysis} in the cytosol converts glucose to pyruvate; the
link reaction and \omterm{def:g10:cellular-respiration:krebs}{Krebs cycle} in the \omterm{def:g10:cellular-respiration:mito}{matrix} release $\mathrm{CO_2}$ and load
electron carriers; the ETC on the inner membrane transfers electrons to
$\mathrm{O_2}$, forming water. Most \omterm{def:g10:cellular-respiration:ATP}{ATP} is made by \omterm{def:g10:cellular-respiration:ETC}{oxidative phosphorylation}
at the ETC/\omterm{def:g10:cellular-respiration:ATP}{ATP} synthase. The final electron acceptor is molecular oxygen.

\textbf{Part II.} \omterm{def:g10:cellular-respiration:lactic}{Lactate} \omterm{def:g10:cellular-respiration:ferment}{fermentation} keeps \omterm{def:g10:cellular-respiration:glycolysis}{glycolysis} running when oxygen
delivery lags. Its \omterm{def:g10:cellular-respiration:ATP}{ATP} yield per glucose is far too small to match aerobic
power output for long, so intense \omterm{def:g10:cellular-respiration:def}{anaerobic} effort cannot be sustained.

\textbf{Part III.} \omterm{def:g10:photosynthesis:def}{Photosynthesis} and respiration share $\mathrm{CO_2}$,
$\mathrm{H_2O}$, $\mathrm{O_2}$ and \omterm{def:g10:molecules-of-life:carbs}{carbohydrate} intermediates (equations are
roughly reverse in outline). A plant in continuous darkness still needs
mitochondria: \omterm{def:g10:photosynthesis:chloroplast}{chloroplasts} do not supply \omterm{def:g10:cellular-respiration:ATP}{ATP} at night, and maintenance,
transport and \omterm{def:g10:scientific-biology:properties}{growth} require respiration of stored organics.
\end{solution}
